\documentclass[12pt]{amsart}
\pagestyle{plain}

\usepackage{amsmath,amssymb}
\usepackage{enumerate}
\usepackage[shortlabels]{enumitem}
\setlist[enumerate]{itemsep=6pt, topsep=6pt, parsep=0pt}
\usepackage{mathtools}
\usepackage{xcolor}
\usepackage[left=0.75in,right=1in,bottom=0.9in]{geometry}
\usepackage[T1]{fontenc}
\usepackage[parfill]{parskip}
\renewcommand{\baselinestretch}{1}

\newcommand{\norm}[1]{\left\|#1\right\|}

% SECTION DIVIDERS
\newcommand{\divider}{
  \vspace{-1em}
  \noindent
  \raisebox{-0.15ex}{\textbullet}%
  \hspace{0.5em}%
  \rule[0.5ex]{0.95\textwidth}{1.0pt}%
  \hspace{0.5em}%
  \raisebox{-0.15ex}{\textbullet}%
  \vspace{0em}
}

\title{ESE 2030 QUIZZAM 4 : SOLUTIONS GUIDE \\ Spring 2026}
\author{prof-g}

\begin{document}
\maketitle

\begin{center}
{\bf NOTE TO STUDENTS}
\end{center}

This solutions guide is organized by course week and topic.
Since there are five versions of this quiz with permuted problem orders \emph{and}
permuted answer choices, the guide does \emph{not} identify correct answers by
letter~(A)--(E).
Instead, the correct answer is \textcolor{blue}{\textbf{highlighted in blue}}
directly within each list of choices.
Match the problem text to the one you received, then read the highlighted
choice, the explanation, and the analysis of each distractor.

Each solution includes:
\begin{itemize}
\item The correct answer (highlighted in blue in the choice list)
\item Detailed mathematical explanation
\item Analysis of every wrong choice, including partial credit designations and common misconceptions
\end{itemize}

\newpage

%%%%%%%%%%%%%%%%%%%%%%%%%%%%%%%%%%%%%%%%%%%%%%%%%%%%%%%%%%%
\section*{WEEK 10: SINGULAR VALUE DECOMPOSITION}
%%%%%%%%%%%%%%%%%%%%%%%%%%%%%%%%%%%%%%%%%%%%%%%%%%%%%%%%%%%

%----------------------------------------------------------
\subsection*{Problem: Frobenius norm from singular values}
%----------------------------------------------------------

A matrix $A \in \mathbb{R}^{m \times n}$ has SVD with singular values
$\sigma_1 \geq \sigma_2 \geq \cdots \geq \sigma_r > 0$.
Which quantity equals $\|A\|_F^2$?

\begin{enumerate}[(A)]
\item $\sigma_1^2$
\item \textcolor{blue}{\textbf{$\sigma_1^2 + \sigma_2^2 + \cdots + \sigma_r^2$}}
\item $\sigma_1 + \sigma_2 + \cdots + \sigma_r$
\item $r \cdot \sigma_1^2$
\item $\sigma_1^2 \cdot \sigma_2^2 \cdots \sigma_r^2$
\end{enumerate}

\textbf{Explanation:}
The Frobenius norm has three equivalent characterizations:
entry-wise $\|A\|_F^2 = \sum_{i,j} a_{ij}^2$;
via trace, $\|A\|_F^2 = \mathrm{tr}(A^TA)$;
and via singular values, $\|A\|_F^2 = \sigma_1^2 + \cdots + \sigma_r^2$.
The equivalence follows from the SVD: $A^TA = V\Sigma^T\Sigma V^T = V\,\mathrm{diag}(\sigma_i^2)\,V^T$,
so $\mathrm{tr}(A^TA) = \sum \sigma_i^2$.

The key conceptual point is that $\|A\|_F$ depends only on the singular values, NOT on
$U$ or $V$. It is therefore invariant under orthogonal multiplication:
$\|QA\|_F = \|AQ\|_F = \|A\|_F$ for any orthogonal $Q$. This invariance is what
makes the Frobenius norm the right object for Eckart--Young--Mirsky.

The contrast with the operator norm is also worth internalizing: $\|A\|_{\mathrm{op}} = \sigma_1$
captures only the maximum stretching direction, whereas the Frobenius norm
captures the ``total stretching energy'' summed across ALL singular directions.
The two norms coincide exactly when $A$ has rank $1$.

\textbf{Trick answer:}
The choice $\sigma_1^2$ is the squared operator norm $\|A\|_{\mathrm{op}}^2$.
Conflating the Frobenius norm with the operator norm is the classic confusion,
since both are standard matrix norms built from the SVD.

\textbf{Trick answer:}
The choice $\sigma_1 + \cdots + \sigma_r$ is the \emph{nuclear norm}
$\|A\|_* = \sum \sigma_i$ --- a different matrix norm used in low-rank
optimization, not the Frobenius norm.

\textbf{Trick answer:}
The choice $r \cdot \sigma_1^2$ equals $\|A\|_F^2$ only when
$\sigma_1 = \sigma_2 = \cdots = \sigma_r$. In general it overestimates,
since smaller singular values contribute less than $\sigma_1^2$ each.

\textbf{Trick answer:}
The product $\sigma_1^2 \cdots \sigma_r^2$ equals $|\det A|^2$ when $A$ is square
of full rank. This is a volume-scaling quantity, not an ``energy'' quantity.

\divider

%----------------------------------------------------------
\subsection*{Problem: Sphere-to-ellipsoid geometry under rank deficiency}
%----------------------------------------------------------

A matrix $A \in \mathbb{R}^{3 \times 3}$ has singular values
$\sigma_1 = 6$, $\sigma_2 = 3$, $\sigma_3 = 0$.
Let $S$ denote the image of the unit sphere
$\{\mathbf{x} \in \mathbb{R}^3 : \|\mathbf{x}\| = 1\}$ under $A$.
Which of the following is TRUE?

\begin{enumerate}[(A)]
\item $S$ is an ellipsoid in $\mathbb{R}^3$ with semi-axes $6$, $3$, and $0$.
\item \textcolor{blue}{\textbf{$S$ is an ellipse in $\mathrm{im}(A)$ with semi-axes $6$ and $3$.}}
\item $S$ is an ellipse in $\ker(A)^\perp$ with semi-axes $6$ and $3$.
\item $S$ is a line segment of length $12$ along $\mathbf{u}_1$.
\item $S$ is an ellipse in $\mathbb{R}^3$ with semi-axes $36$ and $9$.
\end{enumerate}

\textbf{Explanation:}
Decompose $\mathbf{x}$ in the right-singular basis:
$\mathbf{x} = \alpha_1\mathbf{v}_1 + \alpha_2\mathbf{v}_2 + \alpha_3\mathbf{v}_3$
with $\alpha_1^2 + \alpha_2^2 + \alpha_3^2 = 1$. Then
$$A\mathbf{x} = 6\alpha_1\mathbf{u}_1 + 3\alpha_2\mathbf{u}_2,$$
since $\sigma_3 = 0$ kills the $\mathbf{v}_3$ component.
The image $S$ therefore lies in
$\mathrm{span}\{\mathbf{u}_1, \mathbf{u}_2\} = \mathrm{im}(A)$ ---
a $2$-dimensional plane in the codomain ---
and the two surviving stretching factors $6$ and $3$
give the semi-axes.

Two conceptual checks combine in this problem.
First, rank deficiency collapses one dimension, so the image is $2$D, not $3$D.
Second, the image lives in the CODOMAIN (the $\mathbf{u}$-side, spanned by left
singular vectors), NOT in the DOMAIN (the $\mathbf{v}$-side, $\ker(A)^\perp$,
spanned by right singular vectors).

\textbf{Partial credit:}
The choice placing the ellipse in $\ker(A)^\perp$ correctly identifies the image
as a $2$D ellipse with semi-axes $6$ and $3$ --- so the student understands
the rank-deficiency collapse and knows the semi-axes are the $\sigma_i$, not
$\sigma_i^2$. The error is locating the image in the wrong space:
$\ker(A)^\perp \subset \mathbb{R}^3$ is in the DOMAIN, but the image of $A$
lives in the CODOMAIN.

\textbf{Trick answer:}
The choice ``ellipsoid with semi-axis $0$'' refuses to acknowledge that
$\sigma_3 = 0$ collapses one dimension. Geometrically, an ``ellipsoid with
semi-axis $0$'' is just an ellipse in a lower-dimensional plane.

\textbf{Trick answer:}
The line-segment choice over-collapses. Since $\sigma_2 = 3 > 0$ provides a
second independent stretching direction, the image is genuinely $2$D, not $1$D.

\textbf{Trick answer:}
The choice with semi-axes $36$ and $9$ uses $\sigma_i^2$ instead of $\sigma_i$.
The values $36, 9$ are eigenvalues of $A^TA$ --- not the geometric semi-axes
of the image ellipse.

\divider

%----------------------------------------------------------
\subsection*{Problem: Eigen-structure of $A^TA$ and $AA^T$}
%----------------------------------------------------------

Let $A \in \mathbb{R}^{5 \times 3}$ have rank $2$ with singular values
$\sigma_1 \geq \sigma_2 > 0$. Consider the eigen-structure of
$A^TA \in \mathbb{R}^{3 \times 3}$ and $AA^T \in \mathbb{R}^{5 \times 5}$.

Which of the following is TRUE?

\begin{enumerate}[(A)]
\item $A^TA$ and $AA^T$ have the same eigenvalues, counted with multiplicity.
\item $A^TA$ has eigenvalues $\sigma_1, \sigma_2, 0$, and $AA^T$ has eigenvalues
$\sigma_1, \sigma_2, 0, 0, 0$.
\item \textcolor{blue}{\textbf{$A^TA$ and $AA^T$ share the same nonzero eigenvalues; $A^TA$ has one zero eigenvalue and $AA^T$ has three.}}
\item $A^TA$ and $AA^T$ are both invertible, since each is a square matrix.
\item $A^TA$ and $AA^T$ are both positive definite, since each is symmetric with
non-negative eigenvalues.
\end{enumerate}

\textbf{Explanation:}
Both $A^TA$ and $AA^T$ are symmetric and positive semidefinite. Their nonzero
eigenvalues coincide and equal $\sigma_1^2, \sigma_2^2$ (the squared singular
values, two of them since $\mathrm{rank}(A) = 2$). The matrices differ in size,
and so in the number of zero eigenvalues:
\begin{itemize}
\item $A^TA$ is $3\times 3$: eigenvalues $\sigma_1^2,\, \sigma_2^2,\, 0$ (one zero).
\item $AA^T$ is $5\times 5$: eigenvalues $\sigma_1^2,\, \sigma_2^2,\, 0, 0, 0$ (three zeros).
\end{itemize}
In general, the number of zero eigenvalues equals (ambient dimension) $-$ rank.

\textbf{Partial credit:}
The choice listing eigenvalues $\sigma_1, \sigma_2, 0$ and $\sigma_1, \sigma_2, 0, 0, 0$
correctly identifies the structure --- shared nonzero eigenvalues, one zero
for $A^TA$ and three zeros for $AA^T$ --- but uses $\sigma_i$ instead of
$\sigma_i^2$. Captures the harder structural insight while missing the squaring
relationship $\lambda(A^TA) = \sigma^2$.

\textbf{Trick answer:}
``Same eigenvalues counted with multiplicity'' conflates ``same nonzero
eigenvalues'' with ``identical multisets.'' The matrices have different sizes,
so their eigenvalue lists differ in the number of zeros.

\textbf{Trick answer:}
``Square implies invertible'' is false: rank-deficiency forces both products
to be singular regardless of shape.

\textbf{Trick answer:}
Symmetric with non-negative eigenvalues is PSD, not PD; positive definiteness
requires \emph{strictly} positive eigenvalues, ruled out here by rank deficiency.

\divider

%----------------------------------------------------------
\subsection*{Problem: Polar decomposition --- structure and geometric meaning}
%----------------------------------------------------------

Which of the following correctly describes the polar decomposition of an
invertible matrix $A \in \mathbb{R}^{n \times n}$?

\begin{enumerate}[(A)]
\item $A = QR$, where $Q$ is orthogonal and $R$ is upper triangular.
\item \textcolor{blue}{\textbf{$A = QH$, where $Q$ is orthogonal and $H$ is symmetric positive definite, with stretching axes given by the eigenvectors of $A^TA$.}}
\item $A = QH$, where $Q$ is orthogonal and $H$ is symmetric positive definite,
with stretching axes given by the eigenvectors of $A$.
\item $A = QR$, where $Q$ is orthogonal and $R$ is symmetric positive definite.
\item None of the above.
\end{enumerate}

\textbf{Explanation:}
The polar decomposition writes $A = QH$ with $Q$ orthogonal and $H$ symmetric
positive definite. Geometrically, $H$ applies a ``pure stretch'' along orthogonal
axes; $Q$ then rotates (or reflects). Squaring identifies $H$:
$$A^TA = (QH)^T(QH) = H^T Q^T Q H = H^2,$$
so $H = (A^TA)^{1/2}$. The stretching axes of $H$ are therefore the eigenvectors
of $A^TA$ --- equivalently, the right singular vectors of $A$ --- and the
stretching factors are the singular values $\sigma_i$.

A useful analogue: polar coordinates for complex numbers write $z = re^{i\theta}$
as magnitude (stretch) times rotation. The polar decomposition is the matrix
analogue.

\textbf{Partial credit:}
The choice with stretching axes given by ``eigenvectors of $A$'' gets the form
$QH$, the type of each factor (orthogonal, SPD), and the stretch-then-rotate
geometry all correct. Only the wedge is wrong: stretching axes are eigenvectors
of $A^TA$, not of $A$. For non-symmetric $A$, the eigenvectors of $A$ may not
be real or orthogonal, so they cannot serve as the orthogonal stretching axes
of a symmetric matrix.

\textbf{Trick answer:}
``$Q$ orthogonal, $R$ upper triangular'' is the QR decomposition from Week 6.
The shared letter $Q$ tempts students into matching ``orthogonal-times-something''
with the wrong named factorization.

\textbf{Trick answer:}
``$Q$ orthogonal, $R$ symmetric positive definite'' is a hybrid: takes the
structure of polar (orthogonal $\times$ SPD) but uses the QR letter for the
second factor. Catches students who remember ``polar is orthogonal times
something nice'' without keeping the names straight; the actual notation
convention for the second factor is $H$, and a polar-style identity $A = QH$
is not what ``QR decomposition'' refers to.

\textbf{Trick answer:}
``None of the above'' is wrong --- the $A = QH$ description with eigenvectors
of $A^TA$ is the correct polar decomposition.

\divider

%----------------------------------------------------------
\subsection*{Problem: When is an eigendecomposition also an SVD?}
%----------------------------------------------------------

Let $A \in \mathbb{R}^{n \times n}$ be symmetric, with eigendecomposition
$A = Q\Lambda Q^T$.
Under what condition does setting $U = V = Q$ and $\Sigma = \Lambda$
(up to reordering) yield a valid SVD of $A$?

\begin{enumerate}[(A)]
\item Whenever $A$ is symmetric.
\item Whenever $A$ is invertible.
\item \textcolor{blue}{\textbf{Whenever every eigenvalue of $A$ is non-negative.}}
\item Whenever $A$ is positive definite.
\item This eigendecomposition cannot be a valid SVD.
\end{enumerate}

\textbf{Explanation:}
The SVD requires $\Sigma$ to have non-negative diagonal entries.
Taking $U = V = Q$ and $\Sigma = \Lambda$ is a valid SVD precisely when every
$\lambda_i \geq 0$ --- i.e., when $A$ is positive semidefinite.

If some $\lambda_i < 0$, an SVD must use $\sigma_i = |\lambda_i|$ and absorb the
sign by flipping the corresponding column of $U$ or $V$. For example, $A = -I$
has eigenvalues $-1$ but singular values $+1$.

\textbf{Partial credit:}
``$A$ positive definite'' is sufficient but not necessary; PSD is the precise
condition. For instance, $A = \mathrm{diag}(5, 3, 0)$ is PSD (not PD),
and its eigendecomposition is a valid SVD.

\textbf{Trick answer:}
``Whenever $A$ is symmetric'' fails to rule out negative eigenvalues, which
break $\Sigma \geq 0$.

\textbf{Trick answer:}
``Whenever $A$ is invertible'' rules out zero eigenvalues but not negative ones;
$A = -I$ is invertible and symmetric but its eigendecomposition is not an SVD.

\textbf{Trick answer:}
``This cannot be a valid SVD'' is too strong: when $A$ is PSD, the eigen- and
singular value decompositions agree exactly.

\divider

%----------------------------------------------------------
\subsection*{Problem: Reading singular values from an outer-product expansion}
%----------------------------------------------------------

Let $\{\mathbf{u}_1, \mathbf{u}_2\}$ be orthonormal vectors in $\mathbb{R}^4$ and
$\{\mathbf{v}_1, \mathbf{v}_2\}$ be orthonormal vectors in $\mathbb{R}^3$.
A matrix is defined as
$$A = 3\,\mathbf{u}_1\mathbf{v}_1^T + 7\,\mathbf{u}_2\mathbf{v}_2^T.$$
Which of the following is most TRUE?

\begin{enumerate}[(A)]
\item \textcolor{blue}{\textbf{$A$ has rank 2; its singular values are $7$ and $3$.}}
\item $A$ has rank 2; its singular values are $\sqrt{7}$ and $\sqrt{3}$.
\item $A$ has rank 2; its singular values are $49$ and $9$.
\item $A$ has rank 1, since adding two rank-1 terms produces a rank-1 matrix.
\item The singular values of $A$ cannot be determined without knowing the
specific entries of the $\mathbf{u}_i$ and $\mathbf{v}_i$.
\end{enumerate}

\textbf{Explanation:}
An expression $\sum_i c_i\,\mathbf{u}_i\mathbf{v}_i^T$ with orthonormal
$\{\mathbf{u}_i\}$, orthonormal $\{\mathbf{v}_i\}$, and non-negative $c_i$
is already in SVD outer-product form: the $c_i$ ARE the singular values,
and no further computation is needed.

We can verify via $A^TA$. Using $\mathbf{u}_i^T\mathbf{u}_j = \delta_{ij}$:
$$A^TA = 9\,\mathbf{v}_1\mathbf{v}_1^T + 49\,\mathbf{v}_2\mathbf{v}_2^T,$$
whose nonzero eigenvalues are $9$ and $49$. Hence
$\sigma_i = \sqrt{\lambda_i(A^TA)}$ gives $\sigma = 3, 7$ ---
recovering the coefficients.

The rank is $2$ because $\{\mathbf{u}_1, \mathbf{u}_2\}$ are linearly independent
(orthonormal), so the column spaces of the two rank-1 terms intersect trivially.

\textbf{Trick answer:}
The choice $\sqrt{7}, \sqrt{3}$ half-remembers $\sigma_i = \sqrt{\lambda_i(A^TA)}$
but applies the square root to the coefficients $c_i$ themselves rather than to
the eigenvalues of $A^TA$ (which equal $c_i^2$).

\textbf{Trick answer:}
The choice $49, 9$ correctly computes the eigenvalues of $A^TA$ as $9$ and $49$
but forgets the square root step that converts them to singular values.
This is NOT awarded partial credit: the entire point of the problem is reading
singular values directly off an orthonormal expansion without computing $A^TA$
at all. Falling back to $A^TA$ misses the conceptual shortcut.

\textbf{Trick answer:}
The rank-1 claim ignores the orthonormality of the $\mathbf{u}_i$ and
$\mathbf{v}_i$, which guarantees that the column spaces and row spaces of the
two rank-1 terms are mutually orthogonal. The sum has rank $1+1 = 2$.

\textbf{Trick answer:}
``Cannot be determined'' is wrong: orthonormality is precisely the condition
that lets you read singular values directly from the coefficients.
The specific directions affect $U$ and $V$ but not $\Sigma$.

\divider

%----------------------------------------------------------
\subsection*{Problem: Solvability of $A\mathbf{x} = \mathbf{b}$ via fundamental subspaces}
%----------------------------------------------------------

A matrix $A \in \mathbb{R}^{4 \times 3}$ has rank $2$ with SVD $A = U\Sigma V^T$
and left singular vectors $U = [\mathbf{u}_1\;\mathbf{u}_2\;\mathbf{u}_3\;\mathbf{u}_4]$.
Consider $\mathbf{b} = 2\mathbf{u}_1 - \mathbf{u}_2 + 4\mathbf{u}_4$.
Does the system $A\mathbf{x} = \mathbf{b}$ have a solution and, if it does, how many?

\begin{enumerate}[(A)]
\item No solution, because the system is overdetermined (more equations than unknowns).
\item Exactly one solution, because $\mathbf{b}$ has nonzero component in $\mathrm{im}(A)$.
\item Infinitely many solutions, because $\mathrm{rank}(A) = 2 < 3$ means there is one free variable.
\item Exactly one solution, because the pseudoinverse $A^\dagger\mathbf{b}$ is unique.
\item \textcolor{blue}{\textbf{No solution, because $\mathbf{b}$ has a nonzero component in $\mathrm{im}(A)^\perp$.}}
\end{enumerate}

\textbf{Explanation:}
The system $A\mathbf{x} = \mathbf{b}$ is consistent if and only if
$\mathbf{b} \in \mathrm{im}(A) = \mathrm{span}\{\mathbf{u}_1, \mathbf{u}_2\}$.
Decompose $\mathbf{b}$ using the orthogonal direct sum
$\mathbb{R}^4 = \mathrm{im}(A) \oplus \mathrm{im}(A)^\perp$:
$$\mathbf{b} = \underbrace{2\mathbf{u}_1 - \mathbf{u}_2}_{\in\,\mathrm{im}(A)}
+ \underbrace{4\mathbf{u}_4}_{\in\,\mathrm{im}(A)^\perp}.$$
The component $4\mathbf{u}_4 \neq \mathbf{0}$ lies in $\mathrm{im}(A)^\perp$,
so $\mathbf{b}$ does not lie entirely in $\mathrm{im}(A)$. No solution exists.

A useful trichotomy: for $A \in \mathbb{R}^{4 \times 3}$ with rank $2$, either
$\mathbf{b} \notin \mathrm{im}(A)$ and there is no solution, or
$\mathbf{b} \in \mathrm{im}(A)$ and there are infinitely many solutions
(one free variable, since $\dim(\ker A) = 1$). A unique solution is impossible
when $\mathrm{rank}(A) < n$.

\textbf{Trick answer:}
``Overdetermined $\Rightarrow$ no solution'' is wrong: overdetermined systems
CAN have solutions; shape alone never rules them out.

\textbf{Trick answer:}
``Nonzero component in $\mathrm{im}(A) \Rightarrow$ unique solution'' correctly
notes nonzero components along $\mathbf{u}_1, \mathbf{u}_2$ but ignores the
$4\mathbf{u}_4$ component. ALL of $\mathbf{b}$ must lie in $\mathrm{im}(A)$
for the system to be consistent.

\textbf{Trick answer:}
``Wide matrix $\Rightarrow$ infinitely many solutions'' claims an inconsistent
system has infinitely many solutions just because the matrix has more columns
than rank. This conflates uniqueness/nullspace with existence/column space ---
a fatal failure to check the fundamental subspaces.
This is NOT awarded partial credit.

\textbf{Trick answer:}
``$A^\dagger\mathbf{b}$ is unique $\Rightarrow$ unique solution.''
The pseudoinverse $A^\dagger\mathbf{b}$ always exists as the least-squares
solution, but this is not an exact solution to $A\mathbf{x} = \mathbf{b}$
unless $\mathbf{b} \in \mathrm{im}(A)$.

\divider

%----------------------------------------------------------
\subsection*{Problem: Eckart--Young and the optimal rank-$k$ approximation}
%----------------------------------------------------------

Let $A \in \mathbb{R}^{6 \times 4}$ have exactly four nonzero singular values
$\sigma_i$, left singular vectors $\mathbf{u}_i$, and right singular vectors
$\mathbf{v}_i$. Among all rank-$2$ matrices $X$, which minimizes
$\|A - X\|_F$?

\begin{enumerate}[(A)]
\item \textcolor{blue}{\textbf{$\sigma_1\mathbf{u}_1\mathbf{v}_1^T + \sigma_2\mathbf{u}_2\mathbf{v}_2^T$.}}
\item $\sigma_1\mathbf{u}_1\mathbf{v}_1^T + \sigma_4\mathbf{u}_4\mathbf{v}_4^T$.
\item $\sigma_3\mathbf{u}_3\mathbf{v}_3^T + \sigma_4\mathbf{u}_4\mathbf{v}_4^T$.
\item $\sigma_2\mathbf{u}_2\mathbf{v}_2^T + \sigma_3\mathbf{u}_3\mathbf{v}_3^T$.
\item There is not a unique minimizer.
\end{enumerate}

\textbf{Explanation:}
By the Eckart--Young theorem, the truncated SVD
$$A_k = \sum_{i=1}^k \sigma_i\mathbf{u}_i\mathbf{v}_i^T$$
minimizes $\|A - X\|_F$ over ALL rank-$k$ matrices $X$, with error
$\sqrt{\sigma_{k+1}^2 + \cdots + \sigma_r^2}$. With $k = 2$ this gives the
top two terms, with error $\sqrt{\sigma_3^2 + \sigma_4^2}$.

\textbf{Partial credit:}
The choice with $\sigma_1$ and $\sigma_4$ uses the right framework --- a
rank-$2$ truncation built from singular vectors of $A$, with $\sigma_1$
correctly kept --- but is wrong on the second component. The resulting error
is $\sqrt{\sigma_2^2 + \sigma_3^2}$, larger than the true minimum since
$\sigma_2 > \sigma_4$.

\textbf{Trick answer:}
The choice $\sigma_3, \sigma_4$ keeps the SMALLEST two singular components,
giving error $\sqrt{\sigma_1^2 + \sigma_2^2}$ --- the worst among SVD-based
rank-$2$ truncations.

\textbf{Trick answer:}
The choice $\sigma_2, \sigma_3$ keeps middle components and excludes the
largest. The error is $\sqrt{\sigma_1^2 + \sigma_4^2}$, again worse than the
truncated SVD.

\textbf{Trick answer:}
``No unique minimizer'' is false in this regime: Eckart--Young guarantees
uniqueness whenever the singular values at the truncation boundary are
distinct, $\sigma_k > \sigma_{k+1}$.

\divider

%----------------------------------------------------------
\subsection*{Problem: SVD of the pseudoinverse}
%----------------------------------------------------------

Let $A \in \mathbb{R}^{m \times n}$ have rank $r$, with SVD $A = U\Sigma V^T$
and singular values $\sigma_1 > \sigma_2 > \cdots > \sigma_r > 0$.
Suppose $A^\dagger = \tilde U \tilde\Sigma \tilde V^T$ is an SVD of
$A^\dagger$ in standard form (singular values in descending order).
What is $\tilde\sigma_1$ and its corresponding left singular vector
$\tilde{\mathbf{u}}_1$?

\begin{enumerate}[(A)]
\item $\tilde\sigma_1 = \sigma_1$ and $\tilde{\mathbf{u}}_1 = \mathbf{u}_1$.
\item $\tilde\sigma_1 = 1/\sigma_1$ and $\tilde{\mathbf{u}}_1 = \mathbf{u}_1$.
\item $\tilde\sigma_1 = 1/\sigma_1$ and $\tilde{\mathbf{u}}_1 = \mathbf{v}_1$.
\item $\tilde\sigma_1 = 1/\sigma_r$ and $\tilde{\mathbf{u}}_1 = \mathbf{u}_r$.
\item \textcolor{blue}{\textbf{$\tilde\sigma_1 = 1/\sigma_r$ and $\tilde{\mathbf{u}}_1 = \mathbf{v}_r$.}}
\end{enumerate}

\textbf{Explanation:}
From $A = U\Sigma V^T$, the pseudoinverse is $A^\dagger = V\Sigma^\dagger U^T$.
Three coordinated changes occur in passing from the SVD of $A$ to the SVD
of $A^\dagger$:
\begin{enumerate}[(i)]
\item Singular values are inverted: $\sigma_i \mapsto 1/\sigma_i$.
\item The orthogonal factors $U$ and $V$ swap sides: $V$ now contains the
left singular vectors, $U$ the right singular vectors.
\item Descending-order convention forces an index reversal: since
$\sigma_1 > \cdots > \sigma_r$ implies $1/\sigma_r > \cdots > 1/\sigma_1$,
the LARGEST singular value of $A^\dagger$ is $1/\sigma_r$, paired with
the singular vectors that were associated with $\sigma_r$ in the SVD of $A$.
\end{enumerate}
Combining: $\tilde\sigma_1 = 1/\sigma_r$ and $\tilde{\mathbf{u}}_1 = \mathbf{v}_r$.

The geometric picture clarifies the swap. $A$ stretches $\mathbf{v}_i$ by
$\sigma_i$ and outputs along $\mathbf{u}_i$. $A^\dagger$ inverts this, stretching
$\mathbf{u}_i$ by $1/\sigma_i$ and outputting along $\mathbf{v}_i$.
So $\mathbf{u}_i$ becomes an INPUT direction (right singular vector) for
$A^\dagger$, and $\mathbf{v}_i$ becomes an OUTPUT direction (left singular vector).

The five choices form a clean ladder of which transformations the student gets right:
\begin{itemize}
\item $\tilde\sigma_1 = \sigma_1, \tilde{\mathbf{u}}_1 = \mathbf{u}_1$: none.
\item $\tilde\sigma_1 = 1/\sigma_1, \tilde{\mathbf{u}}_1 = \mathbf{u}_1$: reciprocal only.
\item $\tilde\sigma_1 = 1/\sigma_1, \tilde{\mathbf{u}}_1 = \mathbf{v}_1$: reciprocal $+$ swap.
\item $\tilde\sigma_1 = 1/\sigma_r, \tilde{\mathbf{u}}_1 = \mathbf{u}_r$: reciprocal $+$ reversal.
\item $\tilde\sigma_1 = 1/\sigma_r, \tilde{\mathbf{u}}_1 = \mathbf{v}_r$: all three. CORRECT.
\end{itemize}

\textbf{Partial credit:}
The choice $1/\sigma_1, \mathbf{v}_1$ recognizes both that singular values
invert and that $U, V$ swap roles --- the two structural changes from $A$
to $A^\dagger$. It misses only the descending-order convention, and would
be correct if the convention required ascending order.

\textbf{Partial credit:}
The choice $1/\sigma_r, \mathbf{u}_r$ recognizes both that singular values
invert and that the descending convention forces $1/\sigma_r$ to be the
largest. It misses the $U/V$ swap, pairing $\tilde\sigma_1$ with $\mathbf{u}_r$
instead of $\mathbf{v}_r$.

\textbf{Trick answer:}
The choice $\sigma_1, \mathbf{u}_1$ assumes pseudoinverse leaves the SVD untouched.

\textbf{Trick answer:}
The choice $1/\sigma_1, \mathbf{u}_1$ correctly inverts singular values but
assumes the singular vectors are unchanged --- conflates inversion of $\Sigma$
with inversion of $A$.

\divider

\newpage

%%%%%%%%%%%%%%%%%%%%%%%%%%%%%%%%%%%%%%%%%%%%%%%%%%%%%%%%%%%
\section*{WEEK 11: PRINCIPAL COMPONENTS}
%%%%%%%%%%%%%%%%%%%%%%%%%%%%%%%%%%%%%%%%%%%%%%%%%%%%%%%%%%%

%----------------------------------------------------------
\subsection*{Problem: Properties of the covariance matrix}
%----------------------------------------------------------

A centered data matrix $X \in \mathbb{R}^{n \times d}$ has covariance matrix
$C = \frac{1}{n}X^TX$. Which of the following must $C$ satisfy?

\begin{enumerate}[(A)]
\item $C$ is symmetric and positive definite, with $\mathrm{tr}(C)$ equal to the total variance.
\item $C$ is symmetric and positive semidefinite, with $C_{ij}$ equal to the
correlation between variables $i$ and $j$.
\item $C$ is symmetric, with $C_{ij} = 0$ if and only if variables $i$ and $j$
are independent.
\item $C$ is symmetric and positive semidefinite, with eigenvalues equal to $\sigma_i^2(X)$.
\item \textcolor{blue}{\textbf{$C$ is symmetric and positive semidefinite, with $C_{ii}$ equal to the variance of variable $i$.}}
\end{enumerate}

\textbf{Explanation:}
Three properties of $C = X^TX/n$ must hold:
\begin{enumerate}[(1)]
\item Symmetric: $C^T = (X^TX/n)^T = C$.
\item PSD: $\mathbf{z}^TC\mathbf{z} = \|X\mathbf{z}\|^2/n \geq 0$.
(PD only if $X$ has full column rank.)
\item Diagonal entries: $C_{ii} = \tfrac{1}{n}\sum_{k} x_{ki}^2$
(since the data is centered) is the sample variance of variable $i$.
\end{enumerate}

\textbf{Partial credit:}
The choice claiming $C$ is positive definite with trace equal to total variance
is correct on symmetry and the trace identity, but claims PD instead of PSD.
When $n < d$ or when columns are linearly dependent, $C$ has zero eigenvalues
and is only PSD.

\textbf{Trick answer:}
``$C_{ij} =$ correlation'' is wrong: off-diagonal entries are COVARIANCES,
not correlations. The correlation requires normalization
$\rho_{ij} = C_{ij}/\sqrt{C_{ii}C_{jj}}$.

\textbf{Trick answer:}
``$C_{ij} = 0 \Leftrightarrow$ independent'' confuses uncorrelated with
independent. $C_{ij} = 0$ means uncorrelated, but independence is a strictly
stronger condition involving the full joint distribution.

\textbf{Trick answer:}
``Eigenvalues of $C$ equal $\sigma_i^2(X)$'' is off by a factor of $n$:
the eigenvalues of $C$ are $\sigma_i^2(X)/n$, not $\sigma_i^2(X)$.

\divider

%----------------------------------------------------------
\subsection*{Problem: Correlation matrix as cosine similarity}
%----------------------------------------------------------

A centered data matrix $X \in \mathbb{R}^{n \times 3}$ with columns
$\mathbf{x}_1, \mathbf{x}_2, \mathbf{x}_3$ has correlation matrix
$$R = \begin{bmatrix} 1 & -1 & 0 \\ -1 & 1 & 0 \\ 0 & 0 & 1 \end{bmatrix}.$$
Which of the following must be true?

\begin{enumerate}[(A)]
\item $\mathbf{x}_1 \perp \mathbf{x}_2$, and $\mathbf{x}_3$ is orthogonal to both.
\item $\mathbf{x}_2 = -\mathbf{x}_1$, and $\mathbf{x}_3$ is orthogonal to both.
\item \textcolor{blue}{\textbf{$\mathbf{x}_1$ and $\mathbf{x}_2$ are antiparallel; $\mathbf{x}_3$ is orthogonal to both.}}
\item $\mathbf{x}_1, \mathbf{x}_2, \mathbf{x}_3$ are orthonormal in $\mathbb{R}^n$.
\item None of the above.
\end{enumerate}

\textbf{Explanation:}
Correlation between centered columns is the cosine of the angle between them.
Hence $\rho_{12} = -1 \Rightarrow \theta_{12} = 180^\circ$, so
$\mathbf{x}_2 = -c\,\mathbf{x}_1$ for some $c > 0$ (antiparallel; correlation
fixes the angle but not the magnitude). The conditions $\rho_{13} = \rho_{23} = 0$
give $\mathbf{x}_3 \perp \mathbf{x}_1, \mathbf{x}_2$. The diagonal entries
$\rho_{ii} = 1$ are automatic and carry no information about $\|\mathbf{x}_i\|$.

\textbf{Partial credit:}
The choice $\mathbf{x}_2 = -\mathbf{x}_1$ recognizes the antiparallel
relationship but asserts equality with $-\mathbf{x}_1$ rather than
$-c\,\mathbf{x}_1$. Conflates ``perfect negative correlation'' with
``exact negation.''

\textbf{Trick answer:}
``$\mathbf{x}_1 \perp \mathbf{x}_2$'' confuses negative correlation with zero
correlation: $\rho = -1$ is $180^\circ$, not $90^\circ$.

\textbf{Trick answer:}
The orthonormality choice reads $\rho_{ii} = 1$ as $\|\mathbf{x}_i\| = 1$;
the diagonal of any correlation matrix is $1$ by construction. Also misreads
$\rho_{12} = -1$ as orthogonality.

\textbf{Trick answer:}
``None of the above'' is wrong; the antiparallel choice is correct.

\divider

%----------------------------------------------------------
\subsection*{Problem: Sequential definition of the second principal component}
%----------------------------------------------------------

Let $X \in \mathbb{R}^{n \times d}$ be a centered data matrix with covariance
$C = \frac{1}{n}X^TX$, and let $\mathbf{v}_1$ be the first principal component
direction. Which characterization of the second principal component direction
$\mathbf{v}_2$ is correct?

\begin{enumerate}[(A)]
\item The unit vector $\mathbf{v}$ orthogonal to $\mathbf{v}_1$ that maximizes $\mathbf{v}^T C \mathbf{v}$.
\item The unit eigenvector of $C$ whose eigenvalue has the second-largest absolute value.
\item The unit vector $\mathbf{v}$ orthogonal to $\mathbf{v}_1$ that minimizes $\|X\mathbf{v}\|$.
\item The unit vector $\mathbf{v}$ such that $\{\mathbf{v}_1, \mathbf{v}\}$ maximizes
the total variance $\mathbf{v}_1^T C \mathbf{v}_1 + \mathbf{v}^T C \mathbf{v}$.
\item \textcolor{blue}{\textbf{Both (A) and (D).}}
\end{enumerate}

\textbf{Explanation:}
The first listed characterization (``orthogonal to $\mathbf{v}_1$, maximizes
$\mathbf{v}^TC\mathbf{v}$'') is the standard sequential definition of PCA:
after fixing $\mathbf{v}_1$, the second PC maximizes variance among unit
vectors orthogonal to $\mathbf{v}_1$.

The fourth characterization is the equivalent ``joint'' formulation:
choose the orthonormal pair $\{\mathbf{v}_1, \mathbf{v}\}$ maximizing total
captured variance. Since $\mathbf{v}_1^T C \mathbf{v}_1$ is already fixed
(and maximal), this reduces to maximizing $\mathbf{v}^T C \mathbf{v}$ over
unit vectors orthogonal to $\mathbf{v}_1$.

Both characterizations yield the same $\mathbf{v}_2$, so the ``both''
choice is correct.

\textbf{Partial credit:}
The standalone sequential formulation (orthogonal to $\mathbf{v}_1$, maximizes
$\mathbf{v}^TC\mathbf{v}$) is correct as a description of $\mathbf{v}_2$
but misses that the joint formulation also describes the same vector.
Students who pick this understand PCA but didn't recognize the equivalence.

\textbf{Partial credit:}
The standalone joint formulation is correct for the same reason; same
character of partial credit.

\textbf{Trick answer:}
``Eigenvector with second-largest absolute eigenvalue'' is technically
equivalent for PSD covariance matrices (all eigenvalues nonneg, so
``absolute value'' coincides with ``value''), but the phrasing invokes a
distinction that does not apply here. The deeper reason it is wrong is
that PCA \emph{defines} the components via variance maximization;
the eigenvector characterization is a \emph{consequence} of the Spectral Theorem.

\textbf{Trick answer:}
``Minimizes $\|X\mathbf{v}\|$ orthogonal to $\mathbf{v}_1$'' finds the
direction of LEAST variance in the orthogonal complement --- effectively
the LAST principal component restricted to that subspace, not the second.

\divider

%----------------------------------------------------------
\subsection*{Problem: Covariance vs.\ correlation PCA --- when to standardize}
%----------------------------------------------------------

Two centered datasets each have $n$ observations and $d$ variables.
\textbf{Dataset P:} all $d$ variables measure voltage (mV) from various
resistors on a circuit board.
\textbf{Dataset Q:} the $d$ variables measure temperature, humidity,
wind speed, and barometric pressure.
Which choice of PCA method is most appropriate?

\begin{enumerate}[(A)]
\item Correlation PCA for both datasets.
\item Covariance PCA for both datasets.
\item Either method works for either dataset; the two methods produce the same principal components.
\item \textcolor{blue}{\textbf{Covariance PCA for Dataset P; correlation PCA for Dataset Q.}}
\item Correlation PCA for Dataset P; covariance PCA for Dataset Q.
\end{enumerate}

\textbf{Explanation:}
The choice depends on whether differences in variance are physically
meaningful or are scale artifacts.

For Dataset P, all variables share the unit mV. Variance differences reflect
real physical differences --- e.g., one sensor monitors a more variable node
than another. Covariance PCA correctly weights these differences;
standardizing would destroy meaningful information.

For Dataset Q, the variables have incommensurable units. Variance magnitudes
depend on arbitrary unit choices: switching from $^\circ$C to $^\circ$F changes
the temperature variance by a factor of about $3.24$. Covariance PCA would
then let whichever variable has the largest numerical variance dominate.
Correlation PCA standardizes to unit variance, removing the scale artifact.

The rule of thumb: commensurable units $\to$ covariance PCA;
incommensurable units $\to$ correlation PCA.

\textbf{Trick answer:}
``Correlation PCA for both'' standardizes Dataset P unnecessarily, destroying
the meaningful variance differences encoded in the original mV measurements.

\textbf{Trick answer:}
``Covariance PCA for both'' lets the choice of units determine the principal
components for Dataset Q --- a scale-dependent artifact rather than a
data-driven structure.

\textbf{Trick answer:}
``Either method works'' is wrong: the two methods generally produce different
principal components unless all variables already have equal variance
(in which case the methods coincide trivially).

\textbf{Trick answer:}
The reverse choice (correlation for P, covariance for Q) reverses the correct
rule.

\divider

%----------------------------------------------------------
\subsection*{Problem: Interpreting principal component entries}
%----------------------------------------------------------

Correlation PCA on standardized measurements from five environmental sensors
yields the first two PC directions:
$$\mathbf{v}_1 = \frac{1}{\sqrt{5}}\begin{pmatrix} 1 \\ 1 \\ 1 \\ 1 \\ 1 \end{pmatrix},
\qquad \mathbf{v}_2 = \frac{1}{2}\begin{pmatrix} 1 \\ 1 \\ 0 \\ -1 \\ -1 \end{pmatrix}.$$
The five sensors measure (1) $\mathrm{CO}_2$, (2) methane, (3) ozone, (4) wind speed, (5) solar radiation.
Which interpretation is most consistent with these vectors?

\begin{enumerate}[(A)]
\item \textcolor{blue}{\textbf{$\mathbf{v}_1$ is an overall level; $\mathbf{v}_2$ contrasts $\mathrm{CO}_2$/methane against wind/solar, with ozone not participating.}}
\item $\mathbf{v}_1$ is an overall level; in $\mathbf{v}_2$, ozone is the dominant variable.
\item $\mathbf{v}_1$ carries no information because its entries are all equal.
\item $\mathbf{v}_1$ shows all five variables are perfectly correlated; $\mathbf{v}_2$ shows wind/solar are independent of $\mathrm{CO}_2$/methane.
\item $\mathbf{v}_2$ contrasts $\mathrm{CO}_2$ against methane, indicating they are negatively correlated.
\end{enumerate}

\textbf{Explanation:}
Reading PC entries:
\begin{itemize}
\item \emph{Equal positive entries} (in $\mathbf{v}_1$): a ``common mode'' or
overall level. Projection onto $\mathbf{v}_1$ is an equally-weighted average of
standardized variables --- a global factor.
\item \emph{Opposite signs} (in $\mathbf{v}_2$): $\mathrm{CO}_2$ and methane have
positive entries, wind and solar have negative entries. This is a CONTRAST:
when one group is above average, the other tends to be below.
\item \emph{Zero entry}: ozone's zero entry in $\mathbf{v}_2$ means ozone does not
participate in this pattern.
\end{itemize}

\textbf{Trick answer:}
``Ozone is dominant in $\mathbf{v}_2$'' misreads a zero entry. Zero means ``does
not contribute,'' not ``dominant.''

\textbf{Trick answer:}
``$\mathbf{v}_1$ carries no information because entries are all equal'' is the
exact opposite of the truth: a constant vector is highly informative ---
it captures common-factor variation, often the most important pattern in
correlated data.

\textbf{Trick answer:}
``All five variables perfectly correlated; wind/solar independent of CO$_2$/methane''
overstates and misreads. ``Perfectly correlated'' is too strong, and opposite
signs in $\mathbf{v}_2$ indicate negative association in that mode, not
independence.

\textbf{Trick answer:}
``$\mathbf{v}_2$ contrasts $\mathrm{CO}_2$ against methane'' reads opposite
signs where there are none --- $\mathrm{CO}_2$ and methane have the same
sign in $\mathbf{v}_2$, indicating positive association in that mode.

\divider

%----------------------------------------------------------
\subsection*{Problem: Effective dimensionality vs.\ mathematical rank}
%----------------------------------------------------------

A centered data matrix $X \in \mathbb{R}^{500 \times 20}$ has covariance
eigenvalues, in decreasing order:
$$\lambda_1 = 200, \quad \lambda_2 = 180, \quad \lambda_3 = 10,
\quad \lambda_4 = 1.5, \quad \ldots, \quad \lambda_{20} = 0.1.$$
Which statement about reducing this data to its rank-2 PCA approximation
$X_2$ is correct?

\begin{enumerate}[(A)]
\item $X_2 = X$, since the first two principal components span the column space of $X$.
\item \textcolor{blue}{\textbf{$X_2 \neq X$, but the relative error $\|X - X_2\|_F / \|X\|_F$ is small because the discarded eigenvalues are small relative to $\lambda_1$ and $\lambda_2$.}}
\item $X_2 \neq X$, and the relative error $\|X - X_2\|_F / \|X\|_F$ is determined
by the ratio $\lambda_3/\lambda_2$, which measures the size of the spectral gap.
\item $X_2 = X$, since the rank of $X$ equals the number of principal components
retained when a clear elbow is present in the scree plot.
\item $X_2 \neq X$, and rank-2 reduction is inappropriate here because $X$ has
20 nonzero singular values.
\end{enumerate}

\textbf{Explanation:}
The rank-2 PCA approximation $X_2$ is not equal to $X$ --- the matrix $X$ has
rank $20$, so projection onto a $2$-dimensional subspace strictly loses
information. However, the Frobenius error satisfies
$$\|X - X_2\|_F^2 = n \sum_{i \geq 3} \lambda_i \approx n \cdot 12.6,$$
while $\|X\|_F^2 = n \sum_i \lambda_i \approx n \cdot 391.6$. The relative
error is therefore about $\sqrt{12.6/391.6} \approx 18\%$. Small relative error
despite nonzero rank deficiency is exactly what makes rank-2 reduction
appropriate here. Effective dimension $\neq$ mathematical rank.

\textbf{Partial credit:}
The choice ``error governed by $\lambda_3/\lambda_2$'' correctly identifies
$X_2 \neq X$ and references the spectral gap, but the error is governed by
$\sum_{i \geq 3} \lambda_i$, not by the ratio $\lambda_3/\lambda_2$.
A student picking this choice correctly understands the math of the truncation
but conflates the heuristic for \emph{choosing} $k$ (the gap) with the formula
for the \emph{resulting error}.

\textbf{Trick answer:}
``$X_2 = X$ because the first two PCs span the column space'' is mathematically
false: $X_2 = X$ would require the remaining eigenvalues to be exactly zero,
confusing effective dimension with mathematical rank. NOT awarded partial
credit: claiming equality when the tail eigenvalues are strictly nonzero is
a misunderstanding of what a low-rank approximation is.

\textbf{Trick answer:}
``$X_2 = X$ because of the elbow'' conflates mathematical rank with effective
dimension in the opposite direction --- the rank does not equal the number
of retained components just because the scree plot has a clear elbow.

\textbf{Trick answer:}
``Rank-2 reduction is inappropriate here'' correctly identifies $X_2 \neq X$
but draws the wrong conclusion: rank-2 reduction is appropriate \emph{precisely
because} the discarded variance is negligible, not inappropriate because the
rank exceeds 2.

\divider

%----------------------------------------------------------
\subsection*{Problem: SVD-to-PCA dictionary}
%----------------------------------------------------------

A centered data matrix $X \in \mathbb{R}^{100 \times 3}$ has $X^TX$ with
eigenvalues $900$, $400$, and $100$. The data points
$\{\mathbf{x}_1, \ldots, \mathbf{x}_{100}\}$ are the rows of $X$, and
$\mathbf{v}_1, \mathbf{v}_2, \mathbf{v}_3$ are the corresponding principal
component directions. Which of the following is correct?

\begin{enumerate}[(A)]
\item Projecting all $100$ data points onto the first principal component
direction yields scalars whose sample variance is $9$.
\item The image of the unit sphere in $\mathbb{R}^3$ under $X$ is an ellipsoid
in $\mathbb{R}^{100}$ with semi-axis lengths $900$, $400$, and $100$.
\item The total squared distance from the data points to the 2D subspace spanned
by $\mathbf{v}_1$ and $\mathbf{v}_2$ equals $100$.
\item Statements (A) and (B) are both correct.
\item \textcolor{blue}{\textbf{Statements (A) and (C) are both correct.}}
\end{enumerate}

\textbf{Explanation:}
Three SVD-to-PCA conversions are tested simultaneously.

\emph{PC variances.} Projecting $\mathbf{x}_i$ onto $\mathbf{v}_1$ gives the
scalar $\mathbf{x}_i^T\mathbf{v}_1$. The sample variance of these $100$
scalars is the first eigenvalue of the covariance matrix $C = X^TX/n$:
$$\lambda_1(C) = \lambda_1(X^TX)/n = 900/100 = 9. \quad \text{CORRECT.}$$

\emph{Ellipsoid semi-axes.} The image of the unit sphere under $X$ is an
ellipsoid whose semi-axis lengths are the SINGULAR VALUES of $X$, not the
eigenvalues of $X^TX$:
$$\sigma_i = \sqrt{\lambda_i(X^TX)} \;\Rightarrow\; \sigma = 30, 20, 10. \quad \text{INCORRECT.}$$
Also, the ellipsoid lives in $\mathrm{im}(X) \subset \mathbb{R}^{100}$ as a
$3$-dimensional object, not as a generic shape in $\mathbb{R}^{100}$.

\emph{Residual variance.} The total squared distance from the points to the
2D subspace $\mathrm{span}\{\mathbf{v}_1, \mathbf{v}_2\}$ equals the variance
NOT captured by the first two PCs --- i.e., the third eigenvalue of $X^TX$:
$$\sum_{i=1}^{100} \mathrm{dist}\big(\mathbf{x}_i,\, \mathrm{span}\{\mathbf{v}_1,\mathbf{v}_2\}\big)^2
\;=\; \lambda_3(X^TX) \;=\; 100. \quad \text{CORRECT.}$$
Equivalently: $\|X\|_F^2 = \sum_i \lambda_i(X^TX) = 1400$, the rank-2
truncation captures $\sigma_1^2 + \sigma_2^2 = 1300$, and the Pythagorean
identity gives residual $1400 - 1300 = 100$.

So the PC-variance and residual-variance statements are correct,
the ellipsoid statement is wrong, and the answer is the ``both PC-variance
and residual-variance'' choice.

\textbf{Partial credit:}
The standalone PC-variance statement correctly identifies the
$\lambda(X^TX)/n$ scaling but misses that the residual-variance statement
is also correct.

\textbf{Partial credit:}
The standalone residual-variance statement correctly identifies the
deeper of the two facts --- that the residual squared distance under
projection onto the top $k$ PCs equals the sum of the discarded eigenvalues ---
but misses the PC-variance statement.

\textbf{Trick answer:}
The ellipsoid statement uses eigenvalues of $X^TX$ as semi-axis lengths
when the semi-axes are actually the singular values
$\sqrt{900}, \sqrt{400}, \sqrt{100} = 30, 20, 10$.
This is the most common SVD$\leftrightarrow$PCA confusion: forgetting
the square-root step.

\textbf{Trick answer:}
The ``(A) and (B) both correct'' choice combines the correct PC-variance
statement with the trick ellipsoid statement. A student picking this
understands PC variances but conflates the eigenvalues of $X^TX$ with
geometric semi-axis lengths.

\divider

%----------------------------------------------------------
\subsection*{Problem: Eckart--Young as variance maximization (Pythagorean equivalence)}
%----------------------------------------------------------

Let $X \in \mathbb{R}^{n \times d}$ be a centered data matrix and let $X_k$
denote its rank-$k$ truncated SVD. The truncation satisfies the Pythagorean
identity
$$\|X\|_F^2 \;=\; \|X_k\|_F^2 \,+\, \|X - X_k\|_F^2.$$
Eckart--Young says $X_k$ minimizes $\|X - B\|_F$ over rank-$k$ matrices $B$.
PCA says the projection onto the first $k$ principal components maximizes
the variance captured in $k$ dimensions. Why are these two statements equivalent?

\begin{enumerate}[(A)]
\item They aren't equivalent in general; they happen to coincide for centered data.
\item \textcolor{blue}{\textbf{Because $\|X\|_F^2$ is fixed, minimizing $\|X - X_k\|_F^2$ is the same as maximizing $\|X_k\|_F^2$, and the latter is proportional to the captured variance.}}
\item Because the singular vectors $\mathbf{v}_1, \ldots, \mathbf{v}_k$ are
orthonormal, projecting onto them preserves variance exactly.
\item Because the SVD diagonalizes $X$ and PCA diagonalizes $X^TX$, the two
share eigenvectors.
\item They are not equivalent: Eckart--Young minimizes a matrix norm,
while PCA maximizes a scalar quantity.
\end{enumerate}

\textbf{Explanation:}
The Pythagorean identity (given in the stem) is the entire engine. Since
$\|X\|_F^2$ is a constant of the data, the identity says minimizing
$\|X - X_k\|_F^2$ is equivalent to maximizing
$\|X_k\|_F^2 = \sigma_1^2 + \cdots + \sigma_k^2$. Dividing by $n$ gives the
PCA-captured variance, so the two optimizations have the same maximizer.

\textbf{Partial credit:}
``SVD diagonalizes $X$, PCA diagonalizes $X^TX$, shared eigenvectors''
correctly identifies the SVD/PCA shared-eigenvector structure, which is
\emph{why} the truncation lives in the right subspace --- but it doesn't
explain why the \emph{optimization criteria} are equivalent. A student
picking this understands the algebra but hasn't connected the Pythagorean
identity to the dual nature of the optimization.

\textbf{Trick answer:}
``Coincide for centered data'' is wrong: centering affects the
\emph{interpretation} of $\|X_k\|_F^2/n$ as variance, but the
optimization equivalence holds for any $X$.

\textbf{Trick answer:}
``Orthonormal projections preserve variance'' is true but doesn't address
the equivalence of two optimizations. Orthonormal projection preserving
the squared norm of the projected component is a \emph{consequence} of
the Pythagorean identity, not the explanation of the equivalence.

\textbf{Trick answer:}
``Not equivalent because one minimizes a matrix norm, the other maximizes
a scalar'' conflates ``minimizing a Frobenius norm'' with ``the answer being
a matrix.'' Both are scalar optimizations.

\divider

%----------------------------------------------------------
\subsection*{Problem: Johnson--Lindenstrauss target dimension}
%----------------------------------------------------------

The Johnson--Lindenstrauss lemma states that for any set of $n$ points
$\{\mathbf{x}_1, \ldots, \mathbf{x}_n\} \subset \mathbb{R}^d$ and any
$\varepsilon \in (0, 1)$, there exists a linear map
$f : \mathbb{R}^d \to \mathbb{R}^k$ such that for all pairs $i, j$,
$$(1-\varepsilon)\|\mathbf{x}_i - \mathbf{x}_j\|^2 \;\leq\;
\|f(\mathbf{x}_i) - f(\mathbf{x}_j)\|^2 \;\leq\;
(1+\varepsilon)\|\mathbf{x}_i - \mathbf{x}_j\|^2.$$
The lemma's most striking feature is the size of the target dimension $k$
that suffices. On what does $k$ depend?

\begin{enumerate}[(A)]
\item $k$ grows linearly with $d$ and is independent of $n$.
\item $k$ grows logarithmically with $d$ and is independent of $n$.
\item \textcolor{blue}{\textbf{$k$ grows logarithmically with $n$ and is independent of $d$.}}
\item $k$ grows linearly with $n$ and logarithmically with $d$.
\item $k$ depends on the geometric structure of the point set $\{\mathbf{x}_i\}$
and cannot be bounded in terms of $n$ and $d$ alone.
\end{enumerate}

\textbf{Explanation:}
JL gives $k = O(\varepsilon^{-2} \log n)$. The remarkable content of the
lemma is that the target dimension depends only on the \emph{number} of
points and not on the ambient dimension --- $n$ points in $\mathbb{R}^{10^6}$
embed into the same target dimension as $n$ points in $\mathbb{R}^{100}$.
This independence from $d$ is what makes JL useful for high-dimensional data.

\textbf{Partial credit:}
``Logarithmic in $d$, independent of $n$'' correctly identifies the logarithmic
dependence --- which is the second-most-surprising feature of JL after dimension
independence --- but attaches the logarithm to the wrong quantity. A student
picking this has retained that ``JL gives a logarithmic target dimension''
without remembering which input the log is taken of.

\textbf{Trick answer:}
``Linear in $d$, independent of $n$'' says $k$ scales with the ambient
dimension, missing the entire point of the lemma; this would be the naive
guess from someone who has never heard of JL.

\textbf{Trick answer:}
``Linear in $n$, log in $d$'' combines two log-style dependencies but reverses
them; tempting if the student remembers ``log appears somewhere'' but not where.

\textbf{Trick answer:}
``Depends on geometric structure'' would be PCA-like behavior; the power of
JL is precisely that the target dimension is determined by $n$ and
$\varepsilon$ alone, with no inspection of the data itself.

\divider

\newpage

%%%%%%%%%%%%%%%%%%%%%%%%%%%%%%%%%%%%%%%%%%%%%%%%%%%%%%%%%%%
\section*{WEEK 12: NEURAL NETWORKS}
%%%%%%%%%%%%%%%%%%%%%%%%%%%%%%%%%%%%%%%%%%%%%%%%%%%%%%%%%%%

%----------------------------------------------------------
\subsection*{Problem: Why nonlinearity matters --- composition of affine maps}
%----------------------------------------------------------

A network takes input $\mathbf{x} \in \mathbb{R}^3$ and produces output
$\mathbf{y} \in \mathbb{R}^2$ via two blocks $(W_1, \mathbf{b}_1)$ and
$(W_2, \mathbf{b}_2)$ in series, with no activation function between blocks.
Which statement best describes what this network can compute?

\begin{enumerate}[(A)]
\item (approximately) Any continuous function $f:\mathbb{R}^3 \to \mathbb{R}^2$,
since it has a hidden layer.
\item Any quadratic function of $\mathbf{x}$, since composing two blocks
introduces second-order terms.
\item A more limited class of functions than a single-block map.
\item \textcolor{blue}{\textbf{Exactly what a single-block affine map computes; the extra block adds no expressive power.}}
\item It reduces to a single linear map.
\end{enumerate}

\textbf{Explanation:}
Without an activation function between blocks, the network is the composition
of two affine maps. Expanding:
$$\mathbf{y} = W_2(W_1\mathbf{x} + \mathbf{b}_1) + \mathbf{b}_2
= (W_2W_1)\mathbf{x} + (W_2\mathbf{b}_1 + \mathbf{b}_2).$$
The composition is itself an affine map; stacking does not enlarge the
function class. This is exactly why activation functions are essential ---
without a nonlinearity between layers, depth gives no new expressive power.

Note that the composition is generally affine, not linear, since
$W_2\mathbf{b}_1 + \mathbf{b}_2$ is generally nonzero. The distinction matters
for the closely-related distractor.

\textbf{Trick answer:}
``Approximately any continuous function'' is the universal approximation
conclusion --- but UAT requires a nonlinear activation. Without $\sigma$,
the network can only represent affine maps, regardless of hidden width.

\textbf{Trick answer:}
``Any quadratic function'' is a fundamental misunderstanding of composition:
composing linear (or affine) maps gives a linear (or affine) map, not a
quadratic one. No $x_i x_j$ terms appear.

\textbf{Trick answer:}
``A more limited class than a single-block map'' is wrong in the opposite
direction --- the two-block composition is exactly as expressive as a
single block, neither more nor less.

\textbf{Trick answer:}
``Reduces to a single linear map'' drops the affine offset. Biases do not
vanish; they combine into $W_2\mathbf{b}_1 + \mathbf{b}_2$, which is
generally nonzero. The composition is affine, not linear, unless
$\mathbf{b}_1 = \mathbf{b}_2 = \mathbf{0}$.

\divider

%----------------------------------------------------------
\subsection*{Problem: Effect of jointly scaling perceptron weights and bias}
%----------------------------------------------------------

A perceptron with input $\mathbf{x} \in \mathbb{R}^n$, weights
$\mathbf{w} \in \mathbb{R}^n$, bias $b \in \mathbb{R}$, and sigmoid activation
$\sigma$ computes $y = \sigma(\mathbf{w}^T\mathbf{x} + b)$.
Suppose we replace $\mathbf{w}$ with $2\mathbf{w}$ and replace $b$ with $2b$,
but leave $\sigma$ unchanged. Which statement best describes the effect on
the perceptron's output?

\begin{enumerate}[(A)]
\item The set of inputs $\mathbf{x}$ for which $y = 1/2$ is unchanged.
\item The output $y$ doubles for every input $\mathbf{x}$.
\item The decision boundary $\{\mathbf{x}: y = 1/2\}$ shifts by a constant
vector in $\mathbb{R}^n$.
\item The output becomes more sensitive to changes in $\mathbf{x}$ near the
decision boundary, and the boundary itself moves.
\item \textcolor{blue}{\textbf{The output saturates more sharply near the decision boundary, and the boundary itself is unchanged.}}
\end{enumerate}

\textbf{Explanation:}
The decision boundary is the level set
$\{\mathbf{x} : \mathbf{w}^T\mathbf{x} + b = 0\}$ (since $\sigma(0) = 1/2$).
After scaling, the new pre-activation is
$2\mathbf{w}^T\mathbf{x} + 2b = 2(\mathbf{w}^T\mathbf{x} + b)$.
Two consequences:
\begin{enumerate}[(1)]
\item BOUNDARY UNCHANGED. The zero set is identical:
$2(\mathbf{w}^T\mathbf{x} + b) = 0$ iff $\mathbf{w}^T\mathbf{x} + b = 0$.
The hyperplane is the same set of points in $\mathbb{R}^n$.
\item SHARPER SATURATION. Away from the boundary, the pre-activation is twice
as large in magnitude. The sigmoid pushes $y$ closer to $0$ or $1$ for the
same input, and the slope of $y$ along the normal direction at the boundary
itself is doubled.
\end{enumerate}

This is why scaling $(\mathbf{w}, b)$ jointly is a meaningful operation:
it controls ``confidence'' (saturation) without changing ``what is classified
positive'' (boundary). In the limit of large scaling, the sigmoid approaches
a step function with the same boundary --- recovering the hard-threshold
perceptron.

\textbf{Partial credit:}
The choice ``set of $\mathbf{x}$ with $y = 1/2$ is unchanged'' correctly
identifies that the boundary is unchanged --- the harder of the two effects ---
but misses the saturation-sharpening effect on output values away from the
boundary. A student picking this has the geometric invariance right but
doesn't address what changes in the rest of the input space.

\textbf{Trick answer:}
``Output doubles for every input'' treats the sigmoid as if it were linear:
doubling the pre-activation does not double the output. For instance,
$\sigma(0) = 1/2$ remains $1/2$ after scaling, not $1$.

\textbf{Trick answer:}
``Boundary shifts by a constant vector'' is exactly the case the question
avoids by scaling \emph{both} $\mathbf{w}$ and $b$. A shift would result
from scaling $\mathbf{w}$ alone (changing $\{\mathbf{w}^T\mathbf{x} + b = 0\}$
to a parallel hyperplane), but joint scaling preserves the zero set.

\textbf{Trick answer:}
``Sharper sensitivity AND boundary moves'' correctly identifies sharper
sensitivity at the boundary, but incorrectly claims the boundary moves.
A student picking this has internalized that scaling pre-activations sharpens
the sigmoid but hasn't checked that joint scaling preserves the zero set.

\divider

%----------------------------------------------------------
\subsection*{Problem: Single-perceptron representability and linear separability}
%----------------------------------------------------------

A single perceptron with input $\mathbf{x} \in \mathbb{R}^2$ computes
$y = \sigma(\mathbf{w}^T\mathbf{x} + b)$ and classifies an input as
``positive'' when $y \geq 0.5$. Four shaded regions of $\mathbb{R}^2$ are shown:
\textbf{I} (a half-plane below a diagonal line),
\textbf{II} (a disk),
\textbf{III} (a horizontal strip),
\textbf{IV} (a horizontal half-plane $y > 0.3$).
For which subsets can a single perceptron be trained so that its
``positive'' region exactly matches the shaded region?

\begin{enumerate}[(A)]
\item None can be \emph{exactly} matched.
\item Regions I and II only.
\item Regions I, III, and IV only.
\item Regions II and III only.
\item \textcolor{blue}{\textbf{Regions I and IV only.}}
\end{enumerate}

\textbf{Explanation:}
A perceptron's decision boundary in $\mathbb{R}^2$ is the line
$\mathbf{w}^T\mathbf{x} + b = c$, where $c$ is determined by the threshold
$y \geq 0.5$ and the monotonicity of $\sigma$. The ``positive'' region is the
half-plane on one side of this line.
\begin{itemize}
\item Region I (half-plane below the diagonal) is linearly separable. $\checkmark$
\item Region II (disk) requires a curved boundary; not realizable by any
single line. $\times$
\item Region III (horizontal strip) requires the conjunction of two
inequalities (above one line AND below another); a single perceptron
enforces only one. $\times$
\item Region IV (horizontal half-plane $y > 0.3$) is linearly separable,
with weight vector $\mathbf{w} = (0, 1)$. $\checkmark$
\end{itemize}
This limitation --- representability is restricted to half-planes --- is
exactly what hidden layers exist to overcome.

\textbf{Partial credit:}
The choice ``I, III, and IV'' identifies I and IV correctly as half-planes
but mistakenly includes III. A student picking this understands the
half-plane criterion but didn't recognize that a strip requires the
perceptron to enforce two inequalities at once, which is beyond a single
unit's capability.

\textbf{Trick answer:}
``None can be exactly matched'' is too strong. The perceptron is more limited
than the multi-layer networks students will see next, but it can represent
half-planes exactly.

\textbf{Trick answer:}
``Regions I and II only'' groups I and II as ``regions bounded by a single
curve'' --- but linearity, not smoothness, is the relevant criterion.

\textbf{Trick answer:}
``Regions II and III only'' groups II and III as ``bounded'' or ``convex''
regions --- but convexity is necessary, not sufficient. Both convexity AND
being a half-plane are required.

\divider

%----------------------------------------------------------
\subsection*{Problem: Universal Approximation Theorem --- expressive power vs.\ trainability}
%----------------------------------------------------------

The Universal Approximation Theorem is a foundational result about feed-forward
neural networks with at least one hidden layer and a nonlinear activation
function. Which of the following statements most accurately describes what
the theorem claims?

\begin{enumerate}[(A)]
\item A network with one hidden layer can compute any target function exactly,
provided the hidden layer is wide enough.
\item A network with one hidden layer can approximate any target function to
arbitrary accuracy, provided the hidden layer is wide enough.
\item A network with one hidden layer can approximate any target function to
arbitrary accuracy, provided the hidden layer is wide enough and the weights
are found by stochastic gradient descent.
\item \textcolor{blue}{\textbf{A network with one hidden layer can approximate any target function to arbitrary accuracy, provided the hidden layer is wide enough and the right weights and biases exist.}}
\item Approximation to arbitrary accuracy requires at least two hidden layers;
a single hidden layer is fundamentally insufficient.
\end{enumerate}

\textbf{Explanation:}
The Universal Approximation Theorem is an EXISTENCE statement about expressive
power: with enough hidden units and a nonlinear activation, there \emph{exist}
weights and biases such that a one-hidden-layer network approximates the target
function to any desired accuracy. Two important features of the claim:
\begin{enumerate}[(1)]
\item APPROXIMATION, NOT EQUALITY. The network can get arbitrarily close, but
in general cannot compute the target function exactly.
\item EXISTENCE, NOT TRAINABILITY. The theorem says good weights exist; it says
nothing about whether any algorithm (SGD or otherwise) can find them, how
long that would take, or whether training would converge.
\end{enumerate}

\textbf{Partial credit:}
The choice ``approximate any target function, provided the hidden layer is
wide enough'' has the right headline content --- one hidden layer, arbitrary
accuracy, wide enough --- but is missing the existence qualifier. A student
picking this has internalized the theorem's conclusion without distinguishing
it from a constructive guarantee. This is a forgivable elision; the more
careful version is preferred because it makes the existence-vs-construction
distinction explicit.

\textbf{Trick answer:}
``Compute exactly'' is too strong; the theorem gives approximation, not
equality. Confuses ``arbitrarily close'' with ``equal.''

\textbf{Trick answer:}
``Provided the weights are found by SGD'' conflates expressive power with
trainability. The theorem says good weights exist; SGD's ability to find
them is a separate, much harder question.

\textbf{Trick answer:}
``Requires at least two hidden layers'' is wrong: depth is not required;
one hidden layer is provably sufficient. Modern deep networks use depth
for efficiency and inductive bias, not because shallow networks are
fundamentally limited.

\divider

%----------------------------------------------------------
\subsection*{Problem: PCA vs.\ neural network feature maps}
%----------------------------------------------------------

PCA reduces $X \in \mathbb{R}^{n \times d}$ to $k$-dimensional coordinates
via $\mathbf{y}_{\text{PCA}} = V_k^T \mathbf{x}$, where $V_k$ contains the top
$k$ right singular vectors of $X$. A neural network with a hidden layer of
$k$ units produces $\mathbf{y}_{\text{NN}} = \sigma(W_1 \mathbf{x} + \mathbf{b}_1)$,
where $\sigma$ is a nonlinear activation and $W_1, \mathbf{b}_1$ are learned
during training. Two researchers are given the same $X$. Which statement
best describes how their resulting feature maps compare?

\begin{enumerate}[(A)]
\item Both feature maps are determined entirely by $X$, so the two researchers
will obtain features that span the same $k$-dimensional subspace of input space.
\item \textcolor{blue}{\textbf{The PCA feature map is determined by $X$ alone, whereas the neural network's feature map depends on the training objective; the two researchers can obtain genuinely different features from the same $X$.}}
\item The neural network's feature map can always be written as
$\sigma(V_k^T \mathbf{x} + \mathbf{b}_1)$ for some bias $\mathbf{b}_1$, so the
two approaches differ only in the activation function.
\item The PCA feature map is optimal for reconstructing $X$, so any neural
network feature map trained on $X$ will reconstruct $X$ at most as well as
PCA does.
\item The neural network's feature map is a nonlinear function of $\mathbf{x}$,
so it cannot in general be written as a linear projection $V_k^T \mathbf{x}$;
this is the only structural difference between the two approaches.
\end{enumerate}

\textbf{Explanation:}
The genuine structural difference between PCA and neural-network feature
extraction is not just linearity vs.\ nonlinearity --- it is that PCA is
determined by the data alone, while neural network features are determined
by the data \emph{together with} a training objective.

PCA: $V_k$ comes from the SVD of $X$. There is no notion of a task; $V_k$ is
the same regardless of what the features will be used for.

Neural network: $W_1, \mathbf{b}_1$ are learned by minimizing some loss on
training data. Different losses (classification, regression, or different
supervised targets) yield different $W_1$, even with the same $X$.
The features adapt to the task.

This task-dependence is what makes neural-network features powerful for
supervised problems --- they encode whatever about $\mathbf{x}$ is useful for
the goal --- and also what makes them less interpretable: they are no longer
a canonical representation of $X$ itself.

\textbf{Partial credit:}
The choice ``the network's feature map is nonlinear, so cannot be written as
$V_k^T\mathbf{x}$; this is the only structural difference'' correctly
identifies the linear-vs-nonlinear distinction, which is real and important.
The error is the word ``only'': nonlinearity is one structural difference,
but the deeper difference --- visible even when comparing PCA to a
\emph{linear} network feature map --- is the role of the training objective.
A student picking this has internalized the activation-function story but
hasn't recognized that $W_1$ being LEARNED (rather than computed from $X$)
is itself a structural change, independent of $\sigma$.

\textbf{Trick answer:}
``Both determined entirely by $X$'' assumes the network's training objective
is not an additional input. Also subtly conflates ``feature map'' with
``subspace'' --- the PCA features live in a fixed subspace, but the NN
features generally do not even define a subspace, since $\sigma$ is nonlinear.

\textbf{Trick answer:}
``$W_1$ factors through the PCA basis'' claims the network's $W_1$ must be
a linear combination of top singular vectors. This is false: $W_1$ can have
any rows.

\textbf{Trick answer:}
``PCA optimal for reconstruction $\Rightarrow$ NN reconstructs at most as
well'' applies the wrong framing. PCA is optimal for \emph{linear}
reconstruction in Frobenius norm (Eckart--Young), but the comparison to
general nonlinear feature maps is not the right framing --- and a network
trained for a different objective entirely (classification, say) is not even
trying to reconstruct $X$.

\divider

%----------------------------------------------------------
\subsection*{Problem: Backpropagation --- meaning of the error signal $\boldsymbol{\delta}^{(k)}$}
%----------------------------------------------------------

In backpropagation, the algorithm computes at each layer $k$ of a feed-forward
network a vector $\boldsymbol{\delta}^{(k)}$, called the ``error signal''
(or ``error vector'') at that layer. Which statement best describes what
$\boldsymbol{\delta}^{(k)}$ measures, and why backpropagation goes to the
trouble of computing it?

\begin{enumerate}[(A)]
\item The difference between the layer's actual output and the output the layer
should have produced, supplied as a per-layer target during training.
\item \textcolor{blue}{\textbf{The sensitivity of the loss to the layer's pre-activation $\mathbf{z}^{(k)}$: how much $L$ would change in response to a small perturbation of $\mathbf{z}^{(k)}$.}}
\item The prediction error $\mathbf{y} - \mathbf{t}$ scaled down geometrically
with depth, so that earlier layers receive proportionally smaller error signals.
\item The change in weights to be applied at layer $k$ during the gradient
descent update.
\item A local quantity at layer $k$ that depends only on the layer's own weights
and inputs, not on anything downstream.
\end{enumerate}

\textbf{Explanation:}
The error signal at layer $k$ is the gradient of the loss with respect to the
pre-activation at that layer,
$$\boldsymbol{\delta}^{(k)} = \frac{\partial L}{\partial \mathbf{z}^{(k)}},$$
which is by definition a sensitivity: to first order,
$\Delta L \approx \boldsymbol{\delta}^{(k)} \cdot \Delta \mathbf{z}^{(k)}$.
It tells us how much the loss would change if we could nudge the layer's
pre-activation in some direction --- the layer's ``share of responsibility''
for the output error.

Why this is the quantity worth computing: the gradient with respect to any
weight at layer $k$ factors as
$$\frac{\partial L}{\partial W_{ij}^{(k)}} = \delta_i^{(k)} \cdot a_j^{(k-1)}$$
--- the error signal at the unit, times the activation flowing into that
weight. Once $\boldsymbol{\delta}^{(k)}$ is known, every weight and bias
gradient at layer $k$ is one multiplication away. Backpropagation is precisely
the procedure that computes the $\boldsymbol{\delta}^{(k)}$'s recursively from
output to input.

The deeper conceptual content is the absence of intermediate targets.
In supervised training we are given a ground-truth target ONLY at the output.
Hidden layers have no externally-supplied ``correct value.'' Backpropagation's
central insight is that we can DERIVE an effective target signal for each
hidden layer by chain-ruling the output error backward --- and that derived
signal is exactly $\boldsymbol{\delta}^{(k)}$.

\textbf{Trick answer:}
``Difference between actual and target output, supplied per layer'' is the
most seductive wrong answer: it pictures each hidden layer as having its own
target, with $\boldsymbol{\delta}^{(k)}$ playing the role of (actual $-$ target)
layer by layer. There ARE no targets at hidden layers; backprop exists
precisely because we need a way to assign blame to hidden units WITHOUT
knowing what they ``should have'' produced.

\textbf{Trick answer:}
``Prediction error scaled down geometrically with depth'' is wrong: error
signals do not decay geometrically with depth. At each layer they are
transformed by $(W^{(k+1)})^T$ and $\sigma'$, which can shrink, preserve,
or amplify them depending on the weights and activations.
The vanishing gradient phenomenon is one possible behavior of training
dynamics, but it is not part of the definition or purpose of
$\boldsymbol{\delta}^{(k)}$.

\textbf{Trick answer:}
``Change in weights to be applied'' confuses the error signal with the weight
update. The update is
$W^{(k)} \leftarrow W^{(k)} - \eta \, \partial L / \partial W^{(k)}$,
and $\partial L / \partial W^{(k)}$ differs from $\boldsymbol{\delta}^{(k)}$
by a factor of the incoming activation $\mathbf{a}^{(k-1)}$. The error signal
lives in the SPACE of the layer's pre-activations, not the space of its weights.

\textbf{Trick answer:}
``A purely local quantity'' contradicts the chain rule: the loss is computed
only at the output, so the influence of $\mathbf{z}^{(k)}$ on $L$ flows
through every subsequent layer. The recursion
$$\boldsymbol{\delta}^{(k)} = \big((W^{(k+1)})^T \boldsymbol{\delta}^{(k+1)}\big)
\odot \sigma'(\mathbf{z}^{(k)})$$
shows the dependence on layer $k+1$ explicitly. A purely local quantity at
layer $k$ could not possibly carry information about an output-layer loss.

\divider

\end{document}